Nach $t = \tstO$ hat der Stein die Geschwindigkeit
\begin{align}
 v &= \vstF = \gE\cdot\tst = \vst
\end{align}
Da die Geschwindigkeit gleichmässig von null auf $v$ zunimmt, ist die mittlere Geschwindigkeit $v/2$, und der Stein ist um
\begin{align}
 h &= \hstF = \frac{\vst}{2}\cdot\tst = \hst
\end{align}
gefallen. Die potentielle Energie ändert sich um
\begin{align}
 \Delta E_\mathrm{pot} &= \dEpF = -\mst\cdot\gE\cdot\hst = \dEp \approx \result{\dEpP{0}{2}}
\end{align}
und die kinetische Energie um
\begin{align}
 \Delta E_\mathrm{kin} &= \dEkF = \frac{1}{2}\cdot\mst\cdot\left(\vst\right)^2 = \dEk \approx \result{\dEkP{0}{2}}
\end{align}
Die potentielle Energie nimmt um gleich viel ab, wie die kinetische Energie zunimmt: Die mechanische Energie bleibt erhalten.
